\chapter{Conclusions}\label{ch:conclusions}
\thispagestyle{pagestyle}

\section{Summary of Contributions}
\label{sec:contributions}

This thesis presented SHILATE, an open-source pipeline that trains a
deep reinforcement learning agent inside a Unity physics simulation and
deploys the resulting policy onto an Eclipse Leda SDV runtime stack.
The main contributions are as follows.

\textbf{An integrated simulation environment for SDV-targeted RL.} A
Unity 2022 LTS scene was constructed programmatically, comprising a
realistic rear-wheel-drive vehicle physics model, a 21-ray LiDAR-style
sensor, and a circular obstacle course. The environment publishes
observations and consumes control commands through MQTT at 10~Hz, using
the same protocol and topic namespace as the deployment target, so that
no protocol translation is required when moving from training to
deployment.

\textbf{A shaped reward function with quantified milestone bonuses.} A
four-term reward signal was designed and ablated: a continuous progress
reward, a collision penalty, a half-turn milestone bonus, and a
lap-completion bonus. Ablation experiments confirmed that the
intermediate milestone bonus accelerates convergence by approximately
18{,}000 timesteps and that the magnitude of the collision penalty
directly controls the agent's safety margin.

\textbf{A Gymnasium/SB3 training pipeline with live Editor integration.}
The MQTT interface is wrapped in a Gymnasium-compliant environment,
allowing any Stable-Baselines3 algorithm to be substituted without
modification. A Unity Editor window provides one-click training control,
real-time health monitoring across five independently evaluated failure
modes, and live metric graphs.

\textbf{A complete SDV deployment chain.} A containerised
MQTT-to-Kuksa bridge maps simulator telemetry to COVESA VSS paths, and
a Velocitas vehicle application enforces operational limits on the
running RL policy. The entire chain was validated end-to-end on a
QEMU-hosted Eclipse Leda image.

\section{Limitations}
\label{sec:limitations}

\textbf{Single training environment.} SHILATE supports exactly one
Unity Editor instance at real time (\texttt{Time.timeScale = 1}).
Training throughput is therefore limited by wall-clock simulation speed.

\textbf{No sim-to-real gap mitigation.} The sensor model, vehicle
physics, and obstacle geometry are all abstractions. No domain
randomisation or adaptation mechanism is included, so a policy trained
in SHILATE is unlikely to transfer directly to a physical vehicle
without further fine-tuning.

\textbf{Informational-only safety monitor.} The Velocitas application
publishes alerts but does not currently override the agent's control
commands. The deployment chain provides monitoring rather than
enforcement.

\textbf{Limited action space.} The current action space does not include
gear selection or parking-brake inputs, restricting the range of
manoeuvres available to the agent.

\section{Future Work}
\label{sec:future}

\textbf{Multi-environment parallelisation.} Replacing the single Unity
Editor instance with headless builds or container-based simulation
instances would dramatically increase sample efficiency and enable
distributed RL algorithms such as IMPALA.

\textbf{Curriculum learning.} Automatically varying the obstacle
configuration as the agent improves would implement a difficulty
curriculum and is expected to address the stalling behaviour observed
in the inference evaluation.

\textbf{Physical hardware deployment.} Validating the deployment chain
on a physical SDV platform --- such as a CAN-bus-instrumented RC vehicle
running Eclipse Leda on embedded hardware --- would provide direct
evidence of the pipeline's applicability to real automotive systems.

\textbf{Safety monitor upgrade.} Extending the Velocitas application to
publish override control commands when a limit is exceeded would close
the loop from detection to intervention.

\textbf{Richer observation space.} Adding front-facing camera images or
occupancy-grid representations would bring the perception model closer to
production ADAS systems, at the cost of increased training complexity.
